\documentclass[11pt,a4paper]{article}
\usepackage{amsmath,amssymb,amsthm}
\usepackage[margin=2.5cm]{geometry}
\usepackage{hyperref}

\newtheorem{definition}{Definition}[section]
\newtheorem{theorem}[definition]{Theorem}
\newtheorem{proposition}[definition]{Proposition}
\newtheorem{corollary}[definition]{Corollary}
\newtheorem{remark}[definition]{Remark}
\newtheorem{conjecture}[definition]{Conjecture}

\title{Hyperseed Formalization:\\
  Symbolic Ruminations on Non-Symbolic Awareness}
\author{ProtomegaTron (automated formalization)}
\date{2026-09-19 (deepened v2)}

\begin{document}
\maketitle

\begin{abstract}
We formalize Ben Goertzel's ``Symbolic Ruminations on Non-Symbolic
Awareness'' (2022-08-06) within the Hyperseed ontology, incorporating
d-calculus extensions from Hyperseed v2. Non-symbolic awareness is
modeled as direct base experience with fiber suspended, the paradox
of description as fiber self-negation, and meditative practice as
calibration holonomy that improves fiber-base alignment.
\end{abstract}

\tableofcontents

\section{Non-symbolic awareness as base without fiber}

\begin{definition}[Fibered cognition --- claim 1]
\emph{Normal cognition} is fiber-mediated perception: the subject
perceives base space $B$ through the fiber $F$:
\[
  \text{perception}(b) = \sigma(b) \in F(b) \quad \text{(fiber section)}
\]
The fiber structures, categorizes, and filters base experience.
\end{definition}

\begin{definition}[Non-symbolic awareness --- claims 1, 6]
\emph{Non-symbolic awareness} is base perception with fiber suspended:
\[
  \text{perception}_{\text{ns}}(b) = b \quad \text{(direct base contact)}
\]
The fiber $F$ is not destroyed but temporarily set aside --- the section
$\sigma$ is suspended, and the subject contacts the base directly.
\end{definition}

\begin{proposition}[Paradox of description --- claims 2, 5, 7]
\label{prop:paradox}
Describing non-symbolic awareness symbolically is fiber self-negation:
\[
  \sigma_{\text{describe}}(b) = ``\text{there is no } \sigma(b)''
\]
This is a section that asserts the absence of sections --- fiber content
whose meaning is the absence of fiber content. In classical logic this
is contradictory; in paraconsistent logic it is a productive paradox
that inhabits the interzone between symbolic and non-symbolic.
\end{proposition}

\section{Fiber-base calibration}

\begin{definition}[Fiber-base alignment --- claims 3, 8, 9]
The \emph{alignment} between fiber and base measures how accurately the
fiber represents the base:
\[
  \text{align}(F, B) = \text{corr}(\sigma(b), b) \quad
  \text{(correlation between section and base point)}
\]
Misaligned fiber ($\text{align} \ll 1$) produces distorted perception.
Non-symbolic awareness recalibrates alignment by providing direct base
contact that reveals where the fiber has drifted.
\end{definition}

\begin{proposition}[Post-calibration enhancement --- claim 9]
\label{prop:enhance}
After a non-symbolic calibration episode, the restored fiber has
improved alignment:
\[
  \text{align}(F_{t+\epsilon}, B) > \text{align}(F_t, B)
\]
The enhancement comes from correcting drift, not from adding new fiber.
The fiber is the same dimension but better positioned relative to the
base.
\end{proposition}

\section{d-calculus formulation (Hyperseed v2)}

\begin{proposition}[Fiber suspension curvature --- claim 10]
\label{prop:suspension}
The curvature of the fiber suspension transition:
\[
  \|F_{\nabla_{\text{suspend}}}\| = \text{difficulty of entering
    non-symbolic awareness}
\]
\begin{itemize}
  \item Experienced meditators: $\|F_{\nabla_{\text{suspend}}}\| \ll 1$
    (smooth, easy transition).
  \item Beginners: $\|F_{\nabla_{\text{suspend}}}\| \gg 1$ (sharp,
    effortful transition).
\end{itemize}
Practice reduces suspension curvature --- the transition becomes smoother
with repeated traversal (the connection ``wears in'' like a path through
a forest).
\end{proposition}

\begin{proposition}[Base-fiber alignment gradient --- claim 11]
\label{prop:gradient}
The alignment gradient points toward better fiber-base correlation:
\[
  \nabla_{\text{align}} = \nabla \text{corr}(F, B)
\]
Non-symbolic awareness moves the system along this gradient. The
gradient is steepest where fiber-base misalignment is worst ---
non-symbolic awareness provides the most benefit where it's most needed.
\end{proposition}

\begin{proposition}[Calibration holonomy --- claim 12]
\label{prop:holonomy}
A calibration cycle $\gamma$: fiber $\to$ suspend $\to$ base contact
$\to$ restore fiber produces non-trivial holonomy:
\[
  \operatorname{Hol}_\gamma(\Gamma_{\text{calibrate}}) \neq \mathrm{id}
  \implies \text{fiber is changed by calibration}
\]
The holonomy magnitude measures improvement:
\begin{itemize}
  \item $\|\operatorname{Hol}_\gamma - \mathrm{id}\| \gg 1$:
    transformative meditation (large fiber change).
  \item $\|\operatorname{Hol}_\gamma - \mathrm{id}\| \approx 0$:
    ineffective meditation (negligible fiber change).
\end{itemize}
\end{proposition}

\begin{proposition}[Paradox curvature --- claim 13]
\label{prop:boundary}
The curvature at the symbolic/non-symbolic boundary is inherently
non-zero:
\[
  \|F_{\nabla_{\text{boundary}}}\| > 0 \quad \text{(always)}
\]
because the boundary is between fundamentally different modes (fibered
vs.\ unfibered). A paraconsistent connection $\Gamma_{\text{para}}$ at
this boundary enables productive traversal despite non-zero curvature
--- the paradox is navigable even though it can't be eliminated.
\end{proposition}

\begin{proposition}[Non-symbolic gradient --- claim 14]
\label{prop:null}
The direction of non-symbolic awareness in fiber space is the null
direction:
\[
  \nabla_{\text{nonsymbolic}} = -\nabla \|F\| \quad
  \text{(toward zero fiber)}
\]
This is a legitimate direction with well-defined d-calculus geometry:
curvature, holonomy, and parallel transport are all defined along the
approach to zero fiber. The geometry of approaching emptiness is as
rich as the geometry of any other direction --- the d-calculus doesn't
degenerate at $F = 0$, it reveals new structure.
\end{proposition}



%======================================================================
\section{OCO/2 crosswalk (oco/2:2.0.0-alpha.1)}
\emph{Added 2026-10-03. Interpretive mapping onto the Omega Core Ontology
record registry; OCO/2 is an unfrozen alpha candidate.}

\begin{proposition}[Recalibration --- claims 15--18]
Let $a$ be an Assessment before a suspension and $a'$ the Assessment recomputed afterwards from fresh observations only. Recording the pair $(a,a')$ makes the effect of the suspension measurable; claims that the state is indescribable are kept alongside descriptions of it with a contradiction record, so neither is lost and nothing arbitrary follows.
\end{proposition}

\end{document}
